% META: {"id": "TEXP-ARI-EX-027", "chapitre": "TEXP-ARITHMETIQUE", "type_objet": "exercice", "capacites": ["TEXP-ARI-C9"], "capacites_codes": ["C9"], "methodes": ["M9"], "parcours": 1, "competences": ["calculer"], "duree_min": 15, "mode_creation": "ex_nihilo", "sources_inspiration": [], "fichier_tex": "chapitres/TEXP-ARITHMETIQUE/exercices/TEXP-ARI-EX-027.tex", "corrige_tex": "chapitres/TEXP-ARITHMETIQUE/corriges/TEXP-ARI-CO-027.tex", "status": "generated"}
% BEGIN-VERIFY
% import math
% 
% def pgcd(a, b):
%     while b:
%         a, b = b, a % b
%     return a
% 
% def crible(n):
%     marque = [True] * (n + 1)
%     marque[0] = marque[1] = False
%     for p in range(2, int(math.isqrt(n)) + 1):
%         if marque[p]:
%             for multiple in range(p * p, n + 1, p):
%                 marque[multiple] = False
%     return [i for i, ok in enumerate(marque) if ok]
% 
% assert pgcd(84, 120) == 12
% assert pgcd(0, 5) == 5 and pgcd(7, 0) == 7
% for a in range(0, 40):
%     for b in range(0, 40):
%         assert pgcd(a, b) == math.gcd(a, b)
% assert crible(30) == [2, 3, 5, 7, 11, 13, 17, 19, 23, 29]
% END-VERIFY
\begin{exercice}{TEXP-ARI-EX-027}{1}{15}
Écrire, puis exécuter mentalement, les fonctions Python suivantes.
\begin{enumerate}
  \item \texttt{pgcd(a, b)} : algorithme d'Euclide par boucle
        \texttt{while}, renvoyant le PGCD de deux entiers naturels.
  \item \texttt{crible(n)} : crible d'Ératosthène renvoyant la liste des
        nombres premiers inférieurs ou égaux à $n$.
  \item Donner \texttt{pgcd(84, 120)} et \texttt{crible(30)}.
\end{enumerate}
\end{exercice}
